\chapter{Introdução}
\label{cap:introducao}

\section{Motivação e Justificativa}

A sociedade contemporânea apresenta uma dependência crescente em relação a sistemas de software, exigindo que aplicações, especialmente em contextos institucionais e governamentais, sejam desenvolvidas com alto rigor arquitetural e foco em segurança. Plataformas acadêmicas lidam diariamente com dados sensíveis e necessitam de processos eficientes de desenvolvimento e manutenção contínua. 

O Sistema Integrado de Processos Seletivos (SIPROS) da Universidade Federal de Goiás (UFG) foi concebido para centralizar e padronizar editais de seleção. Em seu estado prévio ao semestre letivo de 2026/1, o sistema encontrava-se em uma fase de maturação técnica que demandava evoluções estruturais, tais como: a transição de fluxos simulados (\textit{mockados}) para integrações reais, a consolidação do roteamento e a externalização de configurações, além da adoção de uma camada de autenticação alinhada aos padrões contemporâneos de segurança.

Atuar na resolução desses problemas em um ambiente colaborativo é fundamental. O trabalho em equipe em projetos de software reais traz desafios organizacionais e técnicos que não podem ser simulados em projetos acadêmicos teóricos isolados. A vivência e a aplicação de boas práticas de Engenharia de Software, em um cenário legado e em constante evolução, proporcionam um ganho prático substancial. O relato desta experiência justifica-se por sua valiosa contribuição acadêmica e prática, detalhando como problemas críticos de segurança e consistência de interface podem ser resolvidos. Para estudantes e futuras turmas, este documento serve como um guia de lições aprendidas; para a coordenação da Fábrica de Software, documenta a evolução arquitetural do próprio SIPROS.

\section{Contexto da Experiência}

O presente relato se insere no contexto da disciplina de Prática em Engenharia de Software, ofertada pelo Instituto de Informática da UFG durante o semestre letivo de 2026/1. A Fábrica de Software acadêmica tem como propósito simular o ambiente de mercado e propõe aos discentes a manutenção evolutiva e refatoração de sistemas reais.

O autor do trabalho integrou a "Equipe Dourada", composta por um total de sete discentes, responsável por atuar no desenvolvimento \textit{full-stack} do sistema SIPROS. No que tange à organização do processo de desenvolvimento, a equipe não esteve obrigada a adotar de forma estrita nenhum \textit{framework} ágil específico. Em vez disso, adotou-se uma metodologia de desenvolvimento \textbf{inspirada} nas metodologias ágeis e no \textit{Scrum}, incorporando suas principais características (como as \textit{Sprints}), mas sendo adaptada continuamente para o contexto e necessidades da equipe ao longo do semestre. A coordenação da Fábrica atuou em um papel semelhante ao de um \textit{Product Owner}, definindo as demandas e prioridades, enquanto os discentes possuíam autonomia para o autogerenciamento, divisão de tarefas e aplicação de ferramentas visuais, como quadros \textit{Kanban}.

\section{Objetivo do Trabalho}

Este Trabalho de Conclusão de Curso possui um objetivo principal que delineia a pesquisa, sendo o mesmo desdobrado em metas específicas que orientam a descrição dos resultados obtidos.

\subsection{Objetivo Geral}

O objetivo principal deste trabalho é descrever e analisar a experiência prática de evolução e refatoração \textit{full-stack} do sistema SIPROS no contexto da Fábrica de Software acadêmica da UFG, com ênfase na implementação de mecanismos robustos de segurança, autenticação e na adoção de práticas visando a elevação da qualidade do software.

\subsection{Objetivos Específicos}

Para alcançar o objetivo geral traçado, foram definidos os seguintes objetivos específicos:

\begin{enumerate}
    \item Descrever o ambiente organizacional, a metodologia de trabalho inspirada em métodos ágeis e a dinâmica de colaboração adotada no projeto SIPROS.
    \item Relatar as implementações técnicas voltadas à segurança da informação, destacando a integração com o provedor de identidade do Google via Keycloak, utilizando a extensão PKCE.
    \item Descrever a execução dos processos de Verificação e Validação (V\&V) aplicados para garantir a conformidade arquitetural da interface (\textit{front-end}) com os Casos de Uso estabelecidos.
    \item Apresentar as práticas de Garantia de Qualidade (QA) adotadas pela equipe, evidenciando a estruturação da suíte de testes automatizados e a adoção obrigatória das revisões por pares (\textit{Code Review}).
    \item Refletir criticamente sobre os desafios e gargalos vivenciados, analisando o amadurecimento e a evolução nas \textit{hard skills} e \textit{soft skills} proporcionados por essa vivência de caráter profissional.
\end{enumerate}

\section{Metodologia do Trabalho}

No âmbito metodológico acadêmico, a elaboração do presente trabalho baseia-se na metodologia de \textbf{pesquisa-ação}. Diferentemente de um estudo de caso puramente observacional, a pesquisa-ação caracteriza-se pela intervenção ativa do pesquisador no próprio ambiente para a resolução de um problema real de Engenharia de Software. Neste cenário, o autor atuou ativamente como membro da equipe de desenvolvimento, mapeando inconsistências, projetando soluções de \textit{design} arquitetural e efetuando modificações no código-fonte, participando diretamente da resolução do problema do sistema legado e na consequente implantação de um padrão de segurança seguro e atual.

Para fundamentar as atividades realizadas, a pesquisa se apoiou em literatura técnico-científica reconhecida na área, como o corpo de conhecimento da Engenharia de Software (SWEBOK) para os conceitos de Verificação e Validação, e normas internacionais como a ISO/IEC 25010 para modelagem de qualidade de produto. Complementarmente, consultou-se a documentação técnica pertinente (RFCs do protocolo OAuth 2.0 e manuais oficiais do \textit{Keycloak} e do ecossistema \textit{Angular}) e utilizou-se do acervo de artefatos produzidos internamente pelo projeto (\textit{Issues}, \textit{Pull Requests} e documentação de Casos de Uso). 

Ressalta-se que a metodologia de desenvolvimento de software utilizada no dia a dia da equipe --- inspirada no paradigma ágil --- distingue-se desta metodologia acadêmica de pesquisa-ação e será tratada em maiores detalhes no Capítulo~\ref{cap:relato}.

\section{Organização do Trabalho}

Para organizar o fluxo lógico do documento e guiar a leitura do relato, este trabalho encontra-se estruturado da seguinte forma:

\begin{itemize}
    \item Capítulo 2: Apresenta as bases teóricas, conceituais e tecnológicas requeridas na execução das tarefas, cobrindo tópicos de Qualidade de Software, Verificação e Validação, além das premissas de Segurança em Aplicações Web baseadas em tokens (OAuth 2.0 e PKCE).
    
    \item Capítulo 3: Concentra a descrição detalhada do relato de experiência do discente. Expõe os desafios enfrentados no \textit{onboarding}, as tomadas de decisões arquiteturais na integração com o GCP (\textit{Google Cloud Platform}), o rastreio e mitigação de defeitos e os resultados consolidados na melhoria do SIPROS.
    
    \item Capítulo 4: Corresponde às considerações finais, contendo a reflexão crítica do autor sobre a evolução das competências necessárias ao Engenheiro de Software, um panorama dos Trabalhos Futuros para o projeto e as percepções consolidadas a partir do semestre na Fábrica de Software.
\end{itemize}
